\section{Experiment Configuration Details}
\label{app:experiment-details}

This appendix separates the end-to-end KRR protocol from the protocol used to
isolate the Fourier linear solver.  The saved bundle contains the formal
six-pipeline broad-comparison matrix over \(10^7\)--\(3\times10^8\), explicit
inverse/eigenpair family replays, the earlier EFGP-only regime map, and the
formal matched fixed-system rows.  Each row records the
data set, kernel, \(N\), effective configuration, residual, timing components,
and prediction metrics.  Medians and paired ratios are formed only within one
protocol and never combine timings from different systems.

\subsection{Stage protocols}

\Cref{tab:experiment-stage-protocols} separates the accounting rules.  Stage~1
compares KRR algorithms, so each method pays for its own model
construction.  Stage~2 compares solvers, so every method receives exactly the
same \(A\) and \(\bvec\).

\begin{table}[H]
\centering
\footnotesize
\caption{Roles and accounting of the two formal experimental stages.}
\label{tab:experiment-stage-protocols}
\begin{tabular}{@{}p{0.18\textwidth}p{0.36\textwidth}p{0.36\textwidth}@{}}
\toprule
Quantity & Stage~1: end-to-end KRR & Stage~2: fixed \(A\betavec=\bvec\) \\
\midrule
Primary question
& Accuracy and setup-plus-solve time of complete KRR pipelines from \(10^7\)
to \(3\times10^8\) samples
& Solver acceleration on one identical Fourier system per case \\
Methods
& Primary family replay: EFGP-CG, explicit inverse, localized EigenPro, and
full-grid EigenPro; secondary broad matrix: Nystr\"om KRR, RPCholesky KRR,
EFGP+Jacobi, and the corresponding EFGP rows
& CG; Jacobi-PCG; proposed full-grid, active-inverse, and active-eigenpair PCG \\
System construction
& Each KRR method performs and pays for its own setup
& One shared, hashed \(D,G,\bvec,\lambda\) for all methods \\
Timing unit
& Method-specific setup plus solve; prediction reported separately
& Preconditioner selection/build plus solve; shared Fourier setup separate \\
Selection
& Inverse and eigenpair configurations fixed separately before timing; all
successful raw time and accuracy rows retained
& Configuration frozen before matched timing \\
Repetition
& One warm-up plus five measured repetitions on the common A100
& One warm-up plus five shuffled, paired measured repetitions \\
Arithmetic and check
& Per-method residual or optimization check and test RMSE
& fp64/complex128 and independently recomputed true residual \\
Numerical status
& Complete matrix with explicit resource-limit outcomes
& Complete matched A100 run with verified common system hash \\
\bottomrule
\end{tabular}
\end{table}

The formal fixed-system target is the Synthetic \(N=3\times10^7\) Mat\'ern
case selected by the completed scale campaign.  Its shared construction cannot
be added to or substituted for any end-to-end KRR row.  Likewise, the earlier
regime-map ratios remain a separate historical protocol.

\subsection{Formal Stage 1 end-to-end KRR campaign}

The primary family scale profile uses
\texttt{efgp-standard-cg}, \texttt{efgp-standard-full-eig},
\texttt{ours-binned-inverse}, and
\texttt{ours-binned-active-eig}.  The inverse and active-eigenpair routes
carry separate transferred top-$k$/box/rank fields.  The reporter keeps the
inverse row and minimizes current median total over localized and full-grid
EigenPro for the second family.  It never minimizes across the two families.
The kernel profile applies the same rule separately to SE and Mat\'ern, and
the robustness profile repeats it at every parameter or data-set value.

The secondary broad scale matrix uses
\(N\in\{10^7,3\times10^7,10^8,3\times10^8\}\) on the same declared training
and test artifacts.  Plain EFGP-CG and EFGP+Jacobi run as reference EFGP
pipelines.  The proposed family includes the full-grid spectral route and the
first-order binned default followed by its frozen active inverse/eigenpair
rule.  Nystr\"om
KRR and RPCholesky KRR act on the data-space kernel problem; their landmark or
pivot selection, feature/factor construction, and solve are part of setup and
solve time.  A Fourier-space Nystr\"om or RPCholesky preconditioner is not a
replacement for either data-space row.

For every method, the primary training total is
\(T_{\mathrm{setup}}+T_{\mathrm{solve}}\).  These two components are also
reported separately so that acceleration from faster setup is not attributed
to the solver, or vice versa.  This is the method-owned algorithmic training
total rather than process wall clock: common data-file input, backend creation,
and host-to-device staging are excluded by declaration.  Prediction time is
reported separately.  All successful time--accuracy outcomes remain visible;
reference-relative accuracy labels are descriptive metadata and do not remove
raw timings.  All methods use the same hardware and repeated-timing policy.
Unavailable, resource-limited, or timed-out rows remain explicit rather than
disappearing from a curve.

Exact simple RPCholesky retains its rank-by-\(N\) factor.  At the frozen rank
256 in fp64 this factor exceeds the declared 48-GiB cap from
\(N=3\times10^7\) onward, so a \texttt{resource\_limit} row is the expected
scalability outcome, not a pilot-method substitution.  It counts as a complete
protocol outcome but not as a timed speed comparison.  The completed scale
study then freezes the \(N=3\times10^7\) Synthetic target before Stage~2
timing.

\subsection{Archived regime map and original full-result audit}
\label{app:archived-regime-map}

\begin{sloppypar}
The archived experiment is retained for two purposes: it records the broad
regime map used to choose configurations before the prospective replay, and
it exposes every original candidate rather than only the row printed in the
main text.  The authoritative Mat\'ern files are
\path{diagnostic_tables/paper_table1_candidates.csv} and
\path{diagnostic_tables/paper_table1_selected.csv} inside the tracked
\path{btab_group_a_group_b_group_c_20260703_053504} export bundle.  The SE
rows use the later corrective
\path{btab_group_b_20260704_081848} bundle.  The one-click notebook copies
these into the run directory as
\path{archived_original_full_results.csv} and
\path{archived_original_operation_selected.csv}, and also writes the collapsed
\path{archived_regime_map_two_family.csv}.  Thus the complete row-level result
is part of the saved paper artifact even though it is too large to typeset in
full.

\Cref{tab:archived-selection-hierarchy} states the selection procedure.  It
is important to separate its historical screen from the current reporting
policy.  The old sweep used a training-fit screen to prevent a nonconverged
candidate from becoming the representative timing row.  The present paper
does not use that screen to hide a successful time--accuracy tradeoff: raw
attempts and test RMSE remain in the full-result CSV, and rejected inverse
attempts are marked explicitly when displayed.  The old rule is reported
because it explains exactly how the transferred configurations were obtained.
\end{sloppypar}

\begin{table}[H]
\centering
\footnotesize
\caption{Auditable hierarchy from original candidates to the two-family
tables.  No current-run timing is consulted when configurations are
transferred.}
\label{tab:archived-selection-hierarchy}
\begin{tabular}{@{}p{0.18\textwidth}p{0.73\textwidth}@{}}
\toprule
Step & Rule \\
\midrule
Candidate universe
& Keep every recorded EFGP-CG, active-inverse, localized Box-EigenPro, and
full-grid EigenPro-PCG row, including failed or capped attempts. \\
Historical operation screen
& Within each data set, kernel, and $N$, require the archived convergence
status and training RMSE no larger than $1.10$ times EFGP-CG.  Test RMSE is
not used for selection. \\
Operation winner
& Among rows passing that historical screen, minimize recorded complete
pipeline total separately for Active inverse, Box-EigenPro, and full-grid
EigenPro-PCG. \\
Family collapse
& The inverse operation is the inverse family.  Box-EigenPro and full-grid
EigenPro-PCG are one active-box EigenPro family; choose the smaller operation
winner by total time. \\
Main-text display
& Report EFGP-CG, the inverse-family result, and the active-box EigenPro-family
result together with test RMSE.  Any shown raw row outside the old screen is
marked and excluded from the headline speed claim rather than deleted. \\
Prospective replay
& Transfer only top-$k$, observed box size, and rank.  Rerun one warm-up plus
five measured repetitions; select the faster localized/full-grid eigenpair
row from the current medians.  Archived time never substitutes for replay
time. \\
\bottomrule
\end{tabular}
\end{table}

The archived regime map is interpretable only at the family level.  For SE,
$M=1225$ makes full inversion the usual winner.  For Mat\'ern, Synthetic
selects moderate inverse boxes, whereas Winnebago increasingly selects the
global eigenpair limit as $N$ grows.  These transitions generate the frozen
family-specific configurations used by the revised
\path{family_scale_10m_300m},
\path{family_robustness_at_selected_target}, and
\path{family_kernel_at_30m} notebook profiles.  The original full candidate
CSV remains the source for verifying every alternative box, rank, iteration
count, timing component, and accuracy value.

\subsection{Legacy EFGP scale-map settings}

\Cref{tab:common-experiment-parameters} collects the common parameters and the
two Fourier setup routes. The six training sizes are used for both the
synthetic and USGS tasks. The kernels differ in Fourier dimension because the
same Fourier approximation tolerance is applied to different spectral
densities.

\begin{table}[H]
\centering
\small
\caption{Settings used in the archived exploratory scale map.}
\label{tab:common-experiment-parameters}
\begin{tabular}{@{}ll@{}}
\toprule
Quantity & Value \tabularnewline
\midrule
Training sizes \(N\) & \(10^6,3\times10^6,10^7,3\times10^7,10^8,3\times10^8\) \tabularnewline
Kernel variance & \(1\) \tabularnewline
Lengthscale \(\ell\) & \(0.1\) \tabularnewline
Regularization \(\lambda\) & \(0.1\) \tabularnewline
Fourier approximation tolerance & \(10^{-5}\) \tabularnewline
NUFFT tolerance & \(10^{-10}\) \tabularnewline
Fourier setup routes & direct (baseline); first-order binned (candidates) \tabularnewline
Relative residual tolerance & \(10^{-7}\) \tabularnewline
Maximum CG iterations & \(80000\) (baseline), \(3000\) (preconditioned) \tabularnewline
Fourier dimensions & \(1225\) (SE), \(35721\) (Mat\'ern-\(3/2\)) \tabularnewline
Device & NVIDIA A100-SXM4-40GB \tabularnewline
\bottomrule
\end{tabular}
\end{table}

The relative residual is \(\|b-A\betavec_t\|_2/\|b\|_2\). A run is marked
converged only when this value reaches the tolerance before the iteration
limit. For the synthetic task, the data seeds are 20260421 (training) and 1
(test), and the noise-free test set has \(N/4\) points. The USGS test set is
the fixed held-out portion described in \cref{subsec:exp-setup}. All methods
within a row use the same training and test data. Synthetic training targets
use Gaussian noise with standard deviation \(0.3\).

\subsection{Candidate boxes and ranks}

The requested score budget \(s\) is not generally equal to \(|B|\): the
smallest centered Cartesian box may include additional indices. The actual box
sizes, rather than \(s\), determine memory and runtime. The values represented
in the saved main sweep are summarized in
\cref{tab:candidate-configurations}. Not every listed box--rank combination is
run for every \(N\); the complete Cartesian sweep over the final three
Mat\'ern boxes and five ranks is the one reported in
\cref{tab:box-rank-sweep}.

\begin{table}[H]
\centering
\scriptsize
\setlength{\tabcolsep}{3pt}
\caption{Candidate score budgets, resulting box sizes, and EigenPro ranks in
the saved experiment bundle. Entries in braces are aligned by position.}
\label{tab:candidate-configurations}
\resizebox{\textwidth}{!}{%
\begin{tabular}{llll}
\toprule
Kernel & Operation & Requested \(s\) / observed \(|B|\) & Candidate \(q\) \\
\midrule
SE & inverse
& \(\{512,1024\}/\{625,1225\}\) & -- \\
SE & EigenPro
& \(\{256,512,768\}/\{361,625,961\}\) & \(\{64,128\}\) \\
Mat\'ern & inverse
& \(\{512,728,1024,2048,4096,8192,12288\}\) & -- \\
& & \(\{625,961,1369,2601,5329,10609,15625\}\) & \\
Mat\'ern & EigenPro
& \(\{2048,4096,8192,16384\}/\{2601,5329,10609,21025\}\) & \(\{192,256,320\}\) \\
& & \(\{20480,25720,35721\}/\{25921,32761,35721\}\) & \(\{192,256,320,384,448\}\) \\
\bottomrule
\end{tabular}
}
\end{table}

These sets explain two conventions used in the main text. First, a
configuration is always labeled by the observed \(|B|\), because this is the
dimension on which the operation is performed. Second, \(B=J_m\) is written as
\(P_{\mathrm{inv}}(J_m)\) or \(P_{\mathrm{eig}}(J_m,q)\), without assigning a
new method name.

\subsection{Legacy active-box alternatives on the largest systems}

\Cref{tab:matern-large-method-comparison} reports the alternatives omitted
from the main scalability tables. A full-grid configuration is written with
\(B=J_m\); it is not assigned another method name.

\begin{table}[H]
\centering
\scriptsize
\setlength{\tabcolsep}{4pt}
\caption{Largest Mat\'ern-\(3/2\) systems (\(M=35721\)). ``No candidate met
the screen'' means that no tested inverse converged within 3000 iterations and
passed the training-RMSE check. Speedup compares the recorded complete
pipelines.}
\label{tab:matern-large-method-comparison}
\resizebox{\textwidth}{!}{%
\begin{tabular}{lllrrrrr}
\toprule
Data & \(N\) & Method & Configuration & Iter. & \(T_{\mathrm{tot}}\) & Test RMSE & Speedup \\
\midrule
Synthetic & \(10^8\) & EFGP-CG & -- & 8831 & 10.54 & \(2.95{\times}10^{-3}\) & 1.00 \\
Synthetic & \(10^8\) & \(P_{\mathrm{inv}}(B)\) & \(|B|=5329\) & 187 & 1.76 & \(2.97{\times}10^{-3}\) & 6.00 \\
Synthetic & \(10^8\) & \(P_{\mathrm{eig}}(B,q)\) & \(|B|=35721,\ q=320\) & 495 & 2.04 & \(2.97{\times}10^{-3}\) & 5.16 \\
\addlinespace[1pt]
Synthetic & \(3{\times}10^8\) & EFGP-CG & -- & 11494 & 15.33 & \(1.73{\times}10^{-3}\) & 1.00 \\
Synthetic & \(3{\times}10^8\) & \(P_{\mathrm{inv}}(B)\) & \(|B|=5329\) & 246 & 2.86 & \(1.74{\times}10^{-3}\) & 5.36 \\
Synthetic & \(3{\times}10^8\) & \(P_{\mathrm{eig}}(B,q)\) & \(|B|=35721,\ q=320\) & 643 & 3.34 & \(1.74{\times}10^{-3}\) & 4.58 \\
\addlinespace[1pt]
USGS & \(10^8\) & EFGP-CG & -- & 13539 & 15.46 & \(1.12{\times}10^{-1}\) & 1.00 \\
USGS & \(10^8\) & \(P_{\mathrm{inv}}(B)\) & no candidate met the screen & -- & -- & -- & -- \\
USGS & \(10^8\) & \(P_{\mathrm{eig}}(B,q)\) & \(|B|=35721,\ q=320\) & 264 & 1.60 & \(1.12{\times}10^{-1}\) & 9.65 \\
\addlinespace[1pt]
USGS & \(3{\times}10^8\) & EFGP-CG & -- & 22466 & 27.75 & \(1.11{\times}10^{-1}\) & 1.00 \\
USGS & \(3{\times}10^8\) & \(P_{\mathrm{inv}}(B)\) & no candidate met the screen & -- & -- & -- & -- \\
USGS & \(3{\times}10^8\) & \(P_{\mathrm{eig}}(B,q)\) & \(|B|=35721,\ q=384\) & 320 & 2.90 & \(1.11{\times}10^{-1}\) & 9.56 \\
\bottomrule
\end{tabular}%
}
\end{table}

\subsection{Row selection and timing}

For each data set, kernel, and \(N\), EFGP-CG supplies the baseline training
RMSE. A preconditioned row is eligible if it converges and its training RMSE
is at most \(1.10\) times that baseline. Among eligible rows of a given
operation, the row with the smallest total time is reported. Test RMSE is
copied from that same row; it is not used to choose between candidates. When
no row passes both conditions, the table states that no candidate met the
screen.

Four timing fields are distinguished. Fourier setup includes the spectral
weights, Toeplitz generator, and right-hand side. The baseline rows use direct
cuFINUFFT setup; the preconditioned rows use the first-order binned route. The
saved label for the latter is \texttt{binned\_c1}. The
two routes target the same prescribed Fourier modes but do not hold the
assembled matrix and right-hand side fixed. Preconditioner construction starts after the
setup recorded for its row. PCG time in \cref{tab:timing-breakdown}
excludes construction, while total time includes setup, construction, and
PCG. Prediction is measured separately and is not included in total time.
Consequently, the reported speedup compares complete recorded pipelines. The
matched reruns in \cref{subsec:matched-confirmation} provide the separate
fixed-system solver comparison.

The output bundle also retains the unrounded values used to compute speedup.
Displayed times are rounded to two decimals, so recomputing a ratio from the
printed values can differ slightly from the printed speedup. Failed runs remain
in the raw tables and are excluded only by the stated convergence and fit
rules.

\subsection{Legacy EFGP cost breakdown}
\label{subsec:timing-breakdown}

The alternative configurations for the largest Mat\'ern systems are reported
once, in \cref{tab:matern-large-method-comparison}. On USGS, no tested inverse
passes the screen at \(N=10^8\) or \(3\times10^8\), whereas
\(P_{\mathrm{eig}}(J_m,q)\) gives \(9.65\times\) and \(9.56\times\)
complete-pipeline speedups.

Iteration counts alone do not include the cost of constructing a
preconditioner. \Cref{tab:timing-breakdown} separates the largest Mat\'ern
runs into Fourier setup, preconditioner construction, and PCG after
construction. It also records the two setup routes, so the faster binned setup
is not attributed to preconditioning. The components sum to total time up to
rounding.

\begin{table}[t]
\centering
\footnotesize
\setlength{\tabcolsep}{3.5pt}
\caption{Timing breakdown in seconds for the fastest configuration in each
largest Mat\'ern case. Binned denotes the first-order binned setup route. The
PCG column excludes preconditioner construction.}
\label{tab:timing-breakdown}
\begin{tabular}{llllrrrr}
\toprule
Data & \(N\) & Setup route & Method & Setup & Prec. build & PCG & Total \\
\midrule
Synthetic & \(10^8\) & Direct & EFGP-CG & 1.00 & -- & 9.54 & 10.54 \\
Synthetic & \(10^8\) & Binned & \(P_{\mathrm{inv}}(B)\) & 0.53 & 0.88 & 0.35 & 1.76 \\
Synthetic & \(3{\times}10^8\) & Direct & EFGP-CG & 2.99 & -- & 12.34 & 15.33 \\
Synthetic & \(3{\times}10^8\) & Binned & \(P_{\mathrm{inv}}(B)\) & 1.53 & 0.88 & 0.46 & 2.86 \\
\addlinespace[1pt]
USGS & \(10^8\) & Direct & EFGP-CG & 1.02 & -- & 14.44 & 15.46 \\
USGS & \(10^8\) & Binned & \(P_{\mathrm{eig}}(J_m,320)\) & 0.51 & 0.60 & 0.49 & 1.60 \\
USGS & \(3{\times}10^8\) & Direct & EFGP-CG & 3.43 & -- & 24.32 & 27.75 \\
USGS & \(3{\times}10^8\) & Binned & \(P_{\mathrm{eig}}(J_m,384)\) & 1.55 & 0.74 & 0.61 & 2.90 \\
\bottomrule
\end{tabular}
\end{table}

For the selected inverse, construction is more expensive than the subsequent
PCG solve, but their sum is much smaller than the solve time of the recorded
baseline. The same statement holds for the EigenPro rows. At
\(N=3\times10^8\), the binned setup is about \(54\%\) of each selected
total. Even if construction and PCG took no time, leaving that setup unchanged
could improve the selected totals by only about \(1.9\times\). Further gains
must therefore reduce setup as well as iterations. The table also shows why
candidates are selected by total time rather than iteration count. Dividing
the solve times by their iteration counts shows that PCG costs about
\(1.7\)--\(1.8\) times as much per step, but the saved iteration reduction is
much larger than the one-time build cost expressed in baseline steps.

\subsection{Legacy active-box size and EigenPro rank}
\label{subsec:box-rank-sweep}

\Cref{tab:box-rank-sweep} reports all 15 configurations in the detailed USGS
Mat\'ern sweep at \(N=3\times10^8\). The actual box size is reported after
enclosing the selected Fourier indices; it may exceed the requested top-\(s\)
budget. All converged rows have test RMSE between \(0.11058\) and \(0.11059\),
so this comparison is determined by computation rather than prediction error.

\begin{table}[H]
\centering
\caption{Total time in seconds / PCG iterations for
\(P_{\mathrm{eig}}(B,q)\) on USGS with \(N=3\times10^8\) and \(M=35721\).
The dagger marks a run that did not reach tolerance within 3000 iterations.}
\label{tab:box-rank-sweep}
\begin{tabular}{rccccc}
\toprule
\(|B|\backslash q\) & 192 & 256 & 320 & 384 & 448 \\
\midrule
25921 & \(7.05/3000^\dagger\) & \(8.92/3000^\dagger\) & \(6.99/2783\) & \(6.36/2232\) & \(5.95/1905\) \\
32761 & \(7.12/3000^\dagger\) & \(7.53/3000^\dagger\) & \(6.81/2557\) & \(5.92/1946\) & \(5.64/1639\) \\
35721 & \(3.29/794\) & \(3.01/541\) & \(2.91/408\) & \(\mathbf{2.90/320}\) & \(2.99/272\) \\
\bottomrule
\end{tabular}
\end{table}

For either smaller box, increasing \(q\) reduces the iteration count, yet even
\(q=448\) is slower than the full box with \(q=192\). On the full grid,
increasing \(q\) from 192 to 448 reduces the iterations from 794 to 272, but
total time is lowest at \(q=384\). Thus box size and rank are not
interchangeable, and the rank giving the fewest iterations need not give the
lowest time. The two fastest rows come from single runs, so their small
difference does not establish a universal best rank.

\subsection{Stage 2 fixed-system timing protocol}

For every controlled case, the Fourier weights, Toeplitz generator, right-hand
side, and regularization are constructed once and fingerprinted with SHA-256.
The fingerprint is checked again after all methods finish.  The reported
method set is CG, Jacobi, and the proposed full-grid, active-inverse, and
active-eigenpair family members, as applicable to the grid size.  The method order is reshuffled in
every repetition and the eigensolver seed is fixed.  A timing repetition
reruns score-box selection for every active method, verifies that its frozen
box hash, rule, and effective rank are unchanged, rebuilds the method-specific
preconditioner, starts from zero, solves, synchronizes the GPU at every timing
boundary, and then recomputes the true residual in complex128.  A row is
eligible only if all claimed repetitions satisfy the \(10^{-7}\) tolerance.

The formal target loads the exact \(N=3\times10^7\), noise-0.3 Synthetic
artifact used in Stage~1.  The default active box has
\(|B|=10609\) and rank 320.  Method configurations are transferred from the
prior sweep and frozen before formal timing; no measured repetition tunes a
box, rank, or method.  CG, Jacobi, the default, active EigenPro, the full-grid
spectral family member, and active inverse are all executed.

The common Fourier construction is excluded from the primary solver total
because it is shared by all methods.  The compared total is active-set
selection plus preconditioner construction plus solve; for CG it is the solve
alone.  The formal reporter recomputes this sum and convergence from
repeat-level files.  It verifies equal initial, final, and per-repeat system
fingerprints, the timing-system SHA-256, the embedded canonical configuration
hash, and the component hashes of the weights, Toeplitz generator, and both
right-hand-side representations.  It also checks the common tolerance,
iteration limit, and zero initial vector before emitting
\cref{tab:fixed-system-formal}.  The formal run passes these checks with system
identifier \texttt{12c0a9811168\ldots}; older build-plus-solve artifacts that
omit selection time are not used for that table.

\subsection{Manitowoc provenance, capacity, and reproduction}
\label{subsec:manitowoc-reproduction}

NOAA identifies the independent source as the 641-square-mile Manitowoc work
unit 300411 and links its official USGS EPT endpoint.  It also records that EPT
conversion changes the horizontal CRS to EPSG:3857 without changing NAVD88
GEOID18 elevations \cite{noaa_wi2county_2023}.  The USGS project report gives
the native EPSG:6345 horizontal CRS, EPSG:5703 vertical CRS, and 2023
collection dates \cite{usgs_wi2county_project_report}.  The work-unit report
gives a mean classification-2 density of 11.37 points m\(^{-2}\)
\cite{usgs_wi2county_manitowoc_report}.

The frozen bounding square is centered at \((437100,4884400)\) m in
EPSG:6345 and has side 52000 m.  It contains the entire conforming work-unit
extent.  The preprocessor first uses EPT node bounds only as a coarse filter,
then transforms every candidate point and applies the exact square test.  It
keeps finite class-2 rows, assigns a stable source identifier from the node
rank and point ordinal, and applies SplitMix64 modulo five (residue zero is
test).  An independent hash orders rows within each EPT depth; shallower
additive levels precede deeper levels.  The train/test split is target-blind,
disjoint, and without replacement.  Both coordinates share the fixed 52-km
scale, while a target center and scale fitted on the first one million rows of
the canonical training order are reused at every size and future depth.  The
sidecars verify exact nested 100k, 1M, and 10M prefixes and 2.5 million test
rows for the largest artifact.

The EPT header reports 24,781,638,330 source points across all classes.  The
exact scan through LOD 6 yields 40,787,587 training-bucket and 10,197,897
test-bucket class-2 points.  The official area and mean-ground-density figures
give a planning estimate of 18,876,273,646 ground points, or
15,101,018,917 training-bucket points after the fixed 80/20 split.  This
estimate motivates a streamed route to a 300M/75M artifact, but it is not an
exact deep-level count.  No 300M artifact and no 300M solver run are reported;
both require a deeper exact scan, and the present resident-array solver cannot
load that task on its 4-GB GPU.

Dataset support was not chosen in a blind held-out evaluation.  A first pilot
used a fixed 8-km square in the separate Buffalo work unit
\path{WI_2County_2_B23}.  On its \(N=10^7\) matched system, CG takes 5374
iterations and 10.972 seconds, while the default takes 1717 iterations and
5.863 seconds, a \(1.9109\times\) paired median speedup; the full-grid family member is
faster.  That pilot prompted the move to full-work-unit support before
Manitowoc was screened.  We therefore retain the Buffalo result only here as a
negative geometry pilot and classify Manitowoc as exploratory despite its
independent acquisition.

From the repository root, the checked preprocessing command is:

\begingroup\scriptsize
\begin{verbatim}
PRE=efgp_eigenpro_py.gpu.benchmark_dataset.preprocess_usgs_ept_wi_2county
EPT=https://s3-us-west-2.amazonaws.com/usgs-lidar-public
PREFIX=USGS_EPT_WI_2County_1_B23
python -m "$PRE" \
  --ept-json "$EPT/WI_2County_1_B23/ept.json" \
  --n-train-list 100000,1000000,10000000 \
  --aoi-center-x 437100 --aoi-center-y 4884400 --aoi-side-m 52000 \
  --aoi-crs EPSG:6345 --max-lod-depth 6 \
  --source-project WI_2County_1_B23 \
  --official-mean-ground-density 11.37 \
  --official-work-unit-area-km2 1660.1823787253759 \
  --dataset-stem-prefix "${PREFIX}_full_workunit_ground_elevation"
\end{verbatim}
\endgroup

The matched timing and separate prediction audit are reproduced by:

\begingroup\scriptsize
\begin{verbatim}
CTL=efgp_eigenpro_py.gpu.box_toeplitz_active_block.controlled
STEM=${PREFIX}_full_workunit_ground_elevation_n10000000
OUT=efgp_eigenpro_py/gpu/box_toeplitz_active_block/controlled/outputs
RUN=${OUT}/integrated_manitowoc_n10000000_matern_q256_primary
python -m "$CTL.benchmark" \
  --dataset-stem "$STEM" \
  --n-train 0 --kernel matern --lengthscale 0.1 --nu 1.5 \
  --variance 1 --lambda 0.1 --fourier-eps 1e-5 --nufft-tol 1e-10 \
  --tol 1e-7 --maxiter 6000 --precision fp64 \
  --methods cg,jacobi,default,full-eig \
  --score-tau 1 --box-budget 8192 --inverse-max-size 1024 \
  --rank 256 \
  --eig-tol 1e-3 --eig-maxiter 1280 --warmup-repeats 1 \
  --measured-repeats 5 --method-order-seed 20260823 \
  --eig-seed 0 \
  --nufft-backend none --precompute-chunk-size 1000000 \
  --post-diagnostic-mode cheap --strict-gpu-eig \
  --output-dir "$RUN"

python -m "$CTL.prediction_audit" \
  --config "$RUN/experiment_config.json" \
  --prediction-chunk-size 100000 --warmup-solves 1 \
  --output-dir "$RUN/prediction_audit"
\end{verbatim}
\endgroup

\subsection{Fixed-system bridge to legacy configurations}

\Cref{tab:matched-q128-bridge} revisits the archived rank and box sizes without
choosing the fastest row after timing.  The Mat\'ern columns form the rank-128
configuration bridge and use \(|B|=2601\).  The SE columns are a separate
matched full-inverse control: their score threshold gives \(|B|=841\), and
they also add the archived full grid \(|J_m|=1225\).  Full inversion is not
formed for the Mat\'ern grid with \(M=35721\).  Every displayed method
converges in all five measured repetitions.

\begin{table}[H]
\centering
\scriptsize
\setlength{\tabcolsep}{2.6pt}
\caption{Matched archived-configuration bridge. Each entry is median
iterations / old-runner construction-plus-solve seconds / paired speedup over
CG; active-method score selection is not included.  Entries in different
columns are separate fixed-system cases and are not combined.}
\label{tab:matched-q128-bridge}
\resizebox{\textwidth}{!}{%
\begin{tabular}{lcccc}
\toprule
Method & Synthetic Mat\'ern & USGS Mat\'ern & Synthetic SE & USGS SE \\
\midrule
CG
& \(4246/8.657/1.000\) & \(4905/10.158/1.000\)
& \(2054/3.279/1.000\) & \(1528/2.458/1.000\) \\
Jacobi
& \(1716/4.125/2.111\) & \(8447/20.509/0.497\)
& \(1724/3.262/0.996\) & \(12039/23.711/0.103\) \\
\(P_{\mathrm{inv}}(B)\)
& \(169/2.514/3.464\) & \(1180/5.667/1.793\)
& \(6/0.159/20.677\) & \(12/0.162/15.069\) \\
\(P_{\mathrm{inv}}(J_m)\)
& -- & --
& \(1/0.289/11.356\) & \(1/0.314/7.853\) \\
\(P_{\mathrm{eig}}(B,128)\)
& \(733/2.465/3.522\) & \(1488/4.622/2.198\)
& \(156/0.705/4.643\) & \(12/0.373/6.630\) \\
\(P_{\mathrm{eig}}(J_m,128)\)
& \(704/3.144/2.754\) & \(303/1.846/5.528\)
& \(106/0.586/5.600\) & \(3/0.359/6.803\) \\
\bottomrule
\end{tabular}%
}
\end{table}

For SE, full-grid inversion reduces PCG to one iteration, but its median build
cost is 0.285 seconds on Synthetic and 0.310 seconds on USGS.  The active
inverse builds in 0.144 and 0.135 seconds, respectively, and is therefore
\(1.82\times\) and \(1.94\times\) faster in cold time than the full inverse.
This directly demonstrates why the archived fastest-row rule cannot be
replaced by iteration count alone.

\subsection{Untimed prediction audit}

The controlled prediction audit is deliberately separate from the timing
repetitions.  It rebuilds one system per data set, solves every reported
Stage~2 method against that shared system, rechecks the fingerprint, loads the same held-out
test artifact, and evaluates the prescribed Fourier predictor in chunks of
100000 points.  Coefficients are transferred to the GPU once per method.
Neither the audit solve nor prediction time enters a speedup.  The audit CSV
records the system hash, independently recomputed residual, test RMSE, RMSE
ratio and difference relative to CG, prediction time, test size, and chunk
size.
